\documentclass[11pt,a4paper]{article}
\usepackage[T1]{fontenc}
\usepackage{lmodern}
\usepackage[margin=25mm]{geometry}
\usepackage{amsmath,amssymb,mathtools,bm,mathrsfs}
\usepackage{booktabs,longtable,array,enumitem}
\usepackage{xcolor,microtype}
\usepackage{tocloft}
\usepackage{hyperref}
\usepackage[capitalize,nameinlink]{cleveref}
\hypersetup{colorlinks=true,linkcolor=blue!45!black,urlcolor=blue!45!black,citecolor=blue!45!black,pdftitle={Phase-field fault model: total energy and native FE strain}}
\setlength{\parskip}{4pt}
\setlength{\parindent}{0pt}
\setcounter{tocdepth}{1}
\setlength{\cftbeforesecskip}{2pt}
\setlist{nosep,leftmargin=*}
\allowdisplaybreaks[2]
\newcommand{\B}{\mathcal B}
\newcommand{\F}{\mathcal F}
\newcommand{\Smat}{\bm{\mathcal S}}
\newcommand{\strain}{\bm\varepsilon}
\newcommand{\tauv}{\bm\tau}
\newcommand{\R}{\bm R}
\newcommand{\Ein}{\bm E^{\rm in}}
\newcommand{\Etot}{\bm E^{\rm tot}}
\newcommand{\Eb}{\bm E^{\rm e}_b}
\newcommand{\Ev}{\bm E^{\rm v}_b}
\newcommand{\qcoh}{\tau^{\rm coh}}
\newcommand{\Gc}{\mathcal G_c}
\newcommand{\dd}{\,\mathrm d}
\newcommand{\sym}{\operatorname{sym}}
\newcommand{\dev}{\operatorname{dev}}
\newcommand{\diag}{\operatorname{diag}}
\newcommand{\Ecal}{\mathscr E}
\newcommand{\Dcal}{\mathscr D}
\newcommand{\vect}[1]{\boldsymbol{#1}}
\newcommand{\note}[1]{\par\smallskip\noindent\textbf{Interpretation.} #1\par\smallskip}
\newcommand{\caution}[1]{\par\smallskip\noindent\textbf{Qualification.} #1\par\smallskip}
\newcolumntype{L}[1]{>{\raggedright\arraybackslash}p{#1}}

\title{The current phase-field fault model\\
\large From stored energy to governing equations and the native FE strain algorithm}
\author{Technical reconstruction for Yimin Jin}
\date{Model snapshot: 9 October 2026}

\begin{document}
\maketitle
\begin{abstract}
This note reconstructs the restricted P2 model documented in the native
localization report, the native FE strain report, and the endpoint comparison.
The phase step uses total recoverable bulk and cohesive energy plus fracture
energy, with direct phase irreversibility. It holds the accepted bulk strain
state and surface elastic slip fixed while redistributing their elastic stresses.
The following mechanical step enforces endpoint equilibrium and integrates
the fixed-profile Maxwell equations exactly for interval-constant rates.
A vector field in the native velocity FE space stores the accepted bulk history.
These are three distinct constructions: stored energy, dissipative constitutive
laws, and a sequential time discretization. They are not a single global
minimization problem. Every derivation below specifies what is held fixed.
The current implementation satisfies constrained first-order phase conditions;
the latest report exhibits a saddle point, so local-minimum selection is not
established. The plane-strain subsection records a subsequent source audit and its
consistent 3D embedding. Earlier raw-gradient results remain historical evidence.
\end{abstract}

{\small\setlength{\parskip}{0pt}\tableofcontents}
\clearpage

\section{What is currently selected}
\label{sec:status}

The operative state and step are
\begin{equation}
 \underbrace{(\phi^a,\bm d_R^a,Z^a)}_{\text{accepted state}}
 \ \longrightarrow\ 
 \underbrace{\phi^k,\tauv^*}_{\text{phase relaxation at fixed }\bm d_R^a,Z^a}
 \ \longrightarrow\ 
 \underbrace{(\bm u^k,p^k,V^k)}_{\text{endpoint mechanics}}
 \ \longrightarrow\ 
 (\phi^k,\bm d_R^k,Z^k).
 \label{eq:overview}
\end{equation}
Here $a=k-1$ is the last \emph{accepted} step, not a predictor or a nonlinear
iterate. The quantities $Z$ and $\bm d_R$ are defined below. The governing choices
in this snapshot are:
\begin{enumerate}
 \item $\Ecal=E_b+E_c+F$ drives phase evolution. The old maximum-history field
 $\mathcal H^+$ is diagnostic after initialization; it is not an additional driver.
 \item Phase trials preserve $\R^a$ and $Z^a$, but update $I$, $\chi$, bulk
 stress, and cohesive traction consistently with the trial phase.
 \item The immutable lower bound is $\phi^a$. No healing is admitted.
 \item Mechanics uses the \emph{endpoint} Maxwell stresses. Exact exponential
 constitutive integration does not make the whole coupled step exact in time.
 \item The optional native representation is $\R_h=\mathcal D\bm d_{R,h}$
 in the native velocity space (Q2 in the baseline, Q3 in the diagnostic
 enrichment). Particle stress is a derived sample in this option.
 \item Fracture energy still uses CPDI particle-domain integration. Consistent
 surface normalization and its derivative remain part of the equations.
\end{enumerate}

The latest report~\cite{localization} resolves the former nonlinear stall with
a small exact free-set Newton correction at the existing tolerance. Compatible
side conditions recover a localized band and reproduce the corresponding
coarse one-dimensional discretization. The unresolved issues are spatial stress
accuracy, selection of stable phase branches, and coupled time/energy accuracy.
P3 propagation is not yet qualified.

\subsection{Scope of the derivation}

Unless explicitly labeled otherwise, assume a straight fixed fault in a
two-dimensional, small-strain, incompressible body; fixed mesh and material
coordinates; no stress rotation or transport; constant and identical bulk and
cohesive parameters $G>0$, $\eta>0$; and a fixed normal association region.
The present native option is a fresh serial run with frozen particles and no
restart, geometry evolution, or RSF state evolution. The selected friction
coefficient is $\mu=0.1$ and radiation damping is zero. We retain symbolic
friction and radiation terms to show their location, while clearly distinguishing
them from tested capabilities.

The exact FE history reduction requires more than equality of relaxation
times: its cancellation below uses the same $G$ and the same $\eta$ for both
constituents. It is not yet a prescription for heterogeneous lithospheric
rheology, a rotating/advecting fault, or the full BP3 problem.

\section{Notation and normal localization}
\label{sec:notation}

The domain is $\B$, the sharp reconstructed fault is $\F$, the tangential
coordinate is $\xi$, and the normal coordinate is $\zeta$. For the selected
straight geometry, $\mathrm dA=\mathrm d\xi\,\mathrm d\zeta$.
Surface fields are extended constantly along each associated normal line.
Let
\begin{equation}
 \Smat=\frac{\bm e^n\otimes\bm e^s+\bm e^s\otimes\bm e^n}{2},
 \qquad \operatorname{tr}\Smat=0,\qquad \Smat:\Smat=\frac12.
\end{equation}
In particular, for a horizontal fault, $\tauv:\Smat=\tau_{xy}$.

\begin{longtable}{L{.20\textwidth}L{.55\textwidth}L{.15\textwidth}}
\toprule Symbol & Meaning & Units\\\midrule\endhead
$\bm u$, $\bm d$ & Velocity and physical total displacement & m/s, m\\
$\strain$ & Strain rate; no extra overdot is used on this symbol & s$^{-1}$\\
$\Etot$, $\Ein$, $\R$ & Total strain, committed inelastic strain, and their difference & 1\\
$\bm d_R$ & FE vector encoding $\R$; not total displacement & m\\
$\tauv$, $p$ & Bulk stress variable and Stokes pressure & Pa\\
$s$, $V=\dot s$ & Accumulated scalar slip and slip rate & m, m/s\\
$s^e$, $s^v$ & Elastic and viscous cohesive slip, $s=s^e+s^v$ & m\\
$Z=Gs^e$ & Authoritative surface elastic-slip history & Pa m\\
$\qcoh=Z/I$ & Scalar cohesive traction; called $q$ in reports & Pa\\
$\phi$, $g$, $h$ & Phase, degradation, and $h=g^{-1}-1$ & 1\\
$I$, $\chi=h/I$ & Normal integral and normalized slip distribution & m, m$^{-1}$\\
$\Gc$, $\ell$ & Critical fracture energy and regularization length & J/m$^2$, m\\
$\lambda=\eta/G$ & Maxwell relaxation time & s\\
\bottomrule
\end{longtable}

In the continuum,
\begin{equation}
 \Etot=\mathcal D\bm d,\qquad
 \dot{\Etot}=\strain=\mathcal D\bm u,
 \qquad \nabla\cdot\bm u=0.
 \label{eq:strain-def}
\end{equation}
With an initially trace-free displacement, symmetric and deviatoric symmetric
gradients coincide for the exact incompressible solution. The discrete distinction
is addressed explicitly in \cref{sec:trace}.

\subsection{Degradation and normalization}

The selected AT1/Lorentz functions are
\begin{align}
 \alpha(\phi)&=\phi,\qquad c_0=4\int_0^1\sqrt{\alpha(z)}\dd z=\frac83,\\
 g(\phi)&=\frac{(1-\phi)^2}{(1-\phi)^2+m\phi(1+p\phi)},
 &h(\phi)&=\frac1{g(\phi)}-1
 =\frac{m\phi(1+p\phi)}{(1-\phi)^2}.
 \label{eq:g-h}
\end{align}
The current $p=1$ gives
\begin{equation}
 h'(\phi)=\frac{m(1+3\phi)}{(1-\phi)^3},\qquad
 h''(\phi)=\frac{6m(1+\phi)}{(1-\phi)^4}.
 \label{eq:h-derivatives}
\end{equation}
For the direct continuum normal integral,
\begin{equation}
 I(\xi;\phi)=\int_{\mathcal T_\xi}h(\phi(\xi,\zeta))\dd\zeta,
 \qquad \chi(\xi,\zeta;\phi)=\frac{h(\phi(\xi,\zeta))}{I(\xi;\phi)},
 \qquad \int_{\mathcal T_\xi}\chi\dd\zeta=1.
 \label{eq:normalization}
\end{equation}
The native code instead uses a consistent surface representation $I_h$ of the
direct integrals. The last equality is then an accuracy property, not an exact
pointwise identity for every normal line. The discrete energy must be differentiated
through this representation; see \cref{sec:discrete-phase}.

Finite, positive $I$ and positive evaluated $I_h$ are required. Since
$h\sim(1-\phi)^{-2}$ near saturation, a Q1 normal profile reaching exactly
$\phi=1$ can have a divergent integral. The formal bound $\phi\leq1$ does not
authorize such an inadmissible normalization. No arbitrary cap or denominator
floor is part of the current model. A fully intact normal line with $I=0$ also
requires a separate activation rule; this pre-existing-band model does not
resolve that nucleation issue.

\subsection{Audited three-dimensional plane-strain embedding}
\label{sec:plane-strain}

This subsection is the source-audited correction to the original report's raw
symmetric-gradient convention. Plane strain means zero out-of-plane displacement
and rate, not zero out-of-plane deviatoric stress. For a raw in-plane symmetric
block $A$, define
\[
 \widehat A=\begin{pmatrix}A&0\\0&0\end{pmatrix},\qquad
 \theta=\operatorname{tr}_2 A,\qquad
 D_3(A)=\widehat A-\frac\theta3\bm1_3.
\]
Its stored in-plane block is $D_2(A)=A-\theta\bm1_2/3$; its missing component is
$D_{3,zz}=-\theta/3=-\operatorname{tr}_2 D_2(A)$. This is not the genuinely
2D deviator $A-\theta\bm1_2/2$. For already deviatoric in-plane blocks $T,U$,
use the reduced 3D metric
\[
 \begin{aligned}
 \langle T,U\rangle_3&=T:U+\operatorname{tr}_2T\operatorname{tr}_2U,\\
 \mathrm{II}(T)&=-\tfrac12\langle T,T\rangle_3,
 \qquad J_2(T)=\tfrac12\langle T,T\rangle_3,\\
 \langle D_2(A),D_2(B)\rangle_3&=A:B-\tfrac13\theta_A\theta_B.
 \end{aligned}
\]
The utility \texttt{consistent\_deviator} expects a raw zero-$zz$ plane-strain
block and returns $D_2(A)$. It is not idempotent on its 2D output; applying it
again incorrectly subtracts $\operatorname{tr}_2D_2(A)/3$. The utility
\texttt{consistent\_second\_invariant\_of\_deviatoric\_tensor} expects an
already deviatoric block and reconstructs its $zz$ component. It does not first
remove a raw tensor's trace. Neither helper is being numerically changed.
Here $\mathrm{II}$ is the helper's signed second invariant; $J_2$ is the
positive deviatoric invariant. This is a notation clarification only.

In the remaining equations, symmetric stress/strain symbols and contractions
mean the 3D embedding. When represented by two-dimensional blocks, use
$\langle\cdot,\cdot\rangle_3$ for deviator--deviator contractions. Define
$\mathcal D\bm v=D_3(\sym\nabla_2\bm v)$, $\R_h=\mathcal D\bm d_{R,h}$,
and $\strain_h=\mathcal D\bm u_h$. The in-plane stress, bulk energy and
viscous power reduce to
\begin{align}
 \tauv_2&=2G D_2(\sym\nabla\bm d_R)-2Z\chi\Smat_2,
 &\tau_{zz}&=-\operatorname{tr}_2\tauv_2,\\
 E_b&=\int_\B\frac{\tauv_2:\tauv_2+
                    (\operatorname{tr}_2\tauv_2)^2}{4G}\dd A,
 &\Dcal_b&=\int_\B\frac{\tauv_2:\tauv_2+
                    (\operatorname{tr}_2\tauv_2)^2}{2\eta}\dd A.
\end{align}
The shear slip tensor has zero trace and zero $zz$. Consequently its contraction
with stress needs no extra trace term. Normalization and its adjoint, cohesive
storage/loss, and fracture energy are unchanged. At fixed accepted $\bm d_R,Z$,
$\tau_{zz}$ is independent of $\phi$: the additional bulk trace contribution
has no direct phase derivative. The mechanics and future accepted anchor still
change for nonzero pointwise divergence.

With physical Stokes pressure $p$, the reduced momentum residual and tangent are
\begin{align}
 r_u(\bm w)&=\int_\B[\tauv_2:\sym\nabla\bm w-p\nabla\cdot\bm w]\dd A
              -\ell_{\rm ext}(\bm w),\\
 K_{uu}(\bm w,\delta\bm u)&=2\kappa\int_\B
   [\sym\nabla\bm w:\sym\nabla\delta\bm u
    -\tfrac13(\nabla\cdot\bm w)(\nabla\cdot\delta\bm u)]\dd A.
\end{align}
No extra trace subtraction is applied to the accepted deviatoric history load.
The bulk/slip cross blocks are unchanged because $\operatorname{tr}\Smat=0$.
The FE coefficient update is still
\[
 \bm d_R^k=\beta\bm d_R^a+(\kappa/G)\bm u^k,
 \quad Z^k=\beta Z^a+\kappa V^k,
 \quad\beta=e^{-\Delta t/\lambda},\quad\kappa=\eta(1-\beta).
\]
Applying the linear $D_3\sym\nabla$ proves exactly
$\tauv^k=\beta\tauv^*+2\kappa(\strain^k-\chi^kV^k\Smat)$, including its
$zz$ component, for any $\Delta t/\lambda$. No projection, small-step assumption
or additional accepted state is introduced.

The source audit finds the generic trace-subtraction assembler already present,
but previously registered only for compressible material or prescribed dilation.
The selected incompressible fixture registered neither. Its coefficient is the
material viscosity output $\kappa$, not creep viscosity $\eta$. Register the existing Picard and Newton variants at their respective entry
points (and their matching simplified preconditioners when applicable)
for the optional FE history path, once only. The actual fixture selects Newton;
its trace-term assembler supplies both residual and Jacobian. A Picard-only
registration is insufficient and is retained as a failed diagnostic attempt.
Preserve the default particle path.
The old raw-stress convention could formally absorb its isotropic part through
$p_{\rm raw}=p+2G\operatorname{tr}_2(\sym\nabla\bm d_R)/3$, but no such pressure
redefinition occurs in source, and this field generally lies outside continuous
Q1 pressure. The correction therefore changes the explicit stress/operator;
it does not silently shift the pressure or alter its boundary normalization.

For dynamic-pressure friction, $\sigma_n=\sigma_n^{\rm bg}+p-
\bm n\cdot\tauv_2\bm n$. In the selected prescribed adiabatic-pressure mode,
$\sigma_n=\sigma_n^{\rm bg}+p_{\rm ad}$, with neither Stokes pressure nor
normal deviatoric stress added. These are distinct implemented models; the
selected fixture has $\sigma_n^{\rm bg}=0$, $p_{\rm ad}=10000$ Pa. A nonzero
trace diagnostic does not authorize a compressible physical benchmark or a
change of friction-pressure model.

\subsection{One coupled bulk quadrature}
In the restricted FE option, \texttt{FE bulk normal quadrature points} selects
the Gauss rule on each normal bulk cell (default three for retained comparisons).
Tangential integration splits at actual bulk/fault cell intersections and uses
at least velocity degree plus one Gauss points on each panel, so the polynomial
Stokes products are resolved. The same tensor-product point/weight list is used
for the momentum residual and Jacobian, bulk--fault map and its transpose,
incoming-anchor bulk phase energy/gradient/Hessian, and bulk work/storage/loss
diagnostics. Geometry/association caches are built from this list at fresh
initialization; the setting is fixed during a trajectory. Raising only a
postprocessing rule does not define this updated discrete functional.
Consistent normalization (analytic transverse rays and its projection/adjoint),
sharp-surface quadrature and CPDI fracture integration are unchanged.
The particle-authoritative path retains its native quadrature and defaults.

\section{State variables, stored energy, and dissipation}
\label{sec:energy}

\subsection{The strain decomposition}

Define
\begin{equation}
 \boxed{\R=\Etot-\Ein,\qquad
 \Eb=\R-s^e\chi\Smat,\qquad
 \tauv=2G\Eb=2G\R-2Z\chi\Smat.}
 \label{eq:state-reconstruction}
\end{equation}
The inelastic strain is committed through the rate rule
\begin{equation}
 \dot{\Ein}=\dot{\Ev}+\chi\dot s^v\Smat,
 \qquad \dot s^e=V-\dot s^v.
 \label{eq:inelastic-kinematics}
\end{equation}
Thus it collects bulk viscous strain and the \emph{deposited cohesive viscous
slip strain}. Explicitly,
\begin{equation}
 \Ein(t)=\Ein(0)+\int_0^t
       [\dot{\Ev}(r)+\chi(r)\dot s^v(r)\Smat]\dd r.
 \label{eq:inelastic-memory}
\end{equation}
For an evolving profile this is not $\Ev+\chi(t)s^v(t)\Smat$.
Old inelastic slip is not redistributed whenever the phase changes.
Likewise, adding another $\chi V\Smat$ to \cref{eq:inelastic-kinematics}
would double-count slip. Friction acts on the same $V$ as a dissipative resistance;
it is not a second additive strain channel here.

\subsection{Recoverable and fracture energy}

The cohesive reduction can be obtained from the local rheological choice
\begin{equation}
 \widetilde G=\frac G{h(\phi)},\qquad
 \widetilde\eta=\frac\eta{h(\phi)},\qquad
 \widetilde\lambda=\frac{\widetilde\eta}{\widetilde G}=\lambda.
 \label{eq:local-cohesive-rheology}
\end{equation}
For a normal profile sharing one elastic slip $s^e$, the cohesive elastic strain
is $\chi s^e\Smat$. Its local elastic energy integrates to
\begin{equation}
 \int_{\mathcal T_\xi}\widetilde G|\chi s^e\Smat|^2\dd\zeta
 =\frac{G(s^e)^2}{2I^2}\int_{\mathcal T_\xi}h\dd\zeta
 =\frac{G(s^e)^2}{2I}.
 \label{eq:cohesive-reduction}
\end{equation}
Similarly, integrating $\widetilde\eta|\chi\dot s^v\Smat|^2$ gives
$\eta(\dot s^v)^2/(2I)$. These expressions use their continuous zero-damage
limits where $h=0$; no division by zero is performed in a local intact region.
The local cohesive stress is $2\widetilde G\chi s^e\Smat=2\qcoh\Smat$.
Equation~\eqref{eq:local-cohesive-rheology} is a constitutive choice, not a
consequence of energy balance alone. For projected $I_h$, the surface energy
below is the defined discrete energy; one must not assume that integrating
the local density with unrelated quadrature reproduces it exactly.

The three parts of the state energy are
\begin{align}
 E_b[\R,s^e,\phi]
 &=\int_\B G|\R-s^e\chi(\phi)\Smat|^2\dd A
 =\int_\B\frac{\tauv:\tauv}{4G}\dd A,\label{eq:bulk-energy}\\
 E_c[s^e,\phi]
 &=\int_\F\frac{G(s^e)^2}{2I(\phi)}\dd\xi
 =\int_\F\frac{Z^2}{2G I(\phi)}\dd\xi,\label{eq:coh-energy}\\
 F[\phi]&=\int_\B\frac{\Gc}{c_0\ell}
       [\alpha(\phi)+\ell^2|\nabla\phi|^2]\dd A,\label{eq:fracture-energy}\\
 \boxed{\Ecal=E_b+E_c+F.}\label{eq:total-energy}
\end{align}
In the two-dimensional setting these integrals have units J/m per unit
out-of-plane length; the powers below have units W/m.
The cohesive spring stiffness per unit area is $G/I$, so its traction is
$\partial [G(s^e)^2/(2I)]/\partial s^e=Gs^e/I=\qcoh$.
Pressure is an incompressibility multiplier, not an additional compressive
storage energy in this model. External mechanical loads enter virtual power;
they do not enter the fixed-displacement phase variation.

\subsection{Constitutive choices beyond stored energy}

Stored energy alone does not determine viscosity or friction. The Maxwell
dashpots and interface resistance are additional constitutive choices.
Introduce the viscous dissipation potential
\begin{equation}
 \mathcal R_v(\dot\Ev,\dot s^v;\phi)
  =\int_\B\eta|\dot\Ev|^2\dd A
   +\int_\F\frac{\eta}{2I}(\dot s^v)^2\dd\xi.
 \label{eq:viscous-potential}
\end{equation}
For the positive-slip branch $V\geq0$, define a friction potential at \emph{fixed}
normal stress and state by
\begin{equation}
 \mathcal R_f(V;\sigma_n,\Theta)
   =\int_\F\left[\int_0^V r_f(v;\sigma_n,\Theta)\dd v
                  +\frac12\eta^{\rm dpr}V^2\right]\dd\xi,
 \qquad r_f=\mu(V,\Theta)\sigma_n.
 \label{eq:friction-potential}
\end{equation}
For the tested constant coefficient, the inner integral is $\mu\sigma_n V$.
The physical dissipated power is $r_f V+\eta^{\rm dpr}V^2$, not generally
the value of the potential. For reverse slip one would use the appropriate
signed friction law/subdifferential; that extension is not verified here.

\section{Deriving the continuous mechanical equations}
\label{sec:continuous}

\subsection{Differentiate the energy before imposing constitutive laws}

Let $\tau_s=\tauv:\Smat$ and define its work-conjugate fault average
\begin{equation}
 \mathcal T_\chi[\tauv](\xi)
      =\int_{\mathcal T_\xi}\chi\tau_s\dd\zeta.
 \label{eq:traction-average}
\end{equation}
At a fixed phase, \cref{eq:inelastic-kinematics,eq:state-reconstruction} give
\begin{align}
 \dot\Eb
 &=\strain-\dot\Ev-\chi\dot s^v\Smat
              -\chi(V-\dot s^v)\Smat
  =\strain-\dot\Ev-\chi V\Smat,\\
 \dot E_b+\dot E_c
 &=\int_\B\tauv:\strain\dd A
   -\int_\F\mathcal T_\chi[\tauv]V\dd\xi
   +\int_\F\qcoh V\dd\xi\nonumber\\
 &\quad-\int_\B\tauv:\dot\Ev\dd A
       -\int_\F\qcoh\dot s^v\dd\xi.
 \label{eq:energy-rate-preflow}
\end{align}
If phase evolves, add the variation $D_\phi\Ecal[\dot\phi]$ at fixed
$\Etot,\Ein,s^e$ to the right-hand side, and include $\dot F$ on the left.

The stationarity of this energy rate plus the dissipation potentials, minus
external power and the incompressibility constraint, with respect to the
independent rates produces the following equations. This is a rate-level
derivation at a fixed state, not a joint minimum of $\Ecal$ over all variables.

\subsection{Dashpots and cohesive history}

Variations of $\dot\Ev$ and $\dot s^v$ yield
\begin{equation}
 -\tauv+2\eta\dot\Ev=0,
 \qquad -\qcoh+\frac\eta I\dot s^v=0.
\end{equation}
Consequently,
\begin{equation}
 \boxed{\dot\Ev=\frac{\tauv}{2\eta},\qquad
 \dot s^v=\frac Z\eta,\qquad
 \dot s^e=V-\frac{s^e}{\lambda},\qquad
 \dot Z+\frac Z\lambda=GV,\quad\lambda=\frac\eta G.}
 \label{eq:flow-laws}
\end{equation}
The actual viscous powers are
\begin{equation}
 \Dcal_b=\int_\B\frac{\tauv:\tauv}{2\eta}\dd A,
 \qquad
 \Dcal_c=\int_\F\frac{Z^2}{\eta I}\dd\xi.
 \label{eq:viscous-power}
\end{equation}

\subsection{Bulk equilibrium and slip balance}

Let $\ell_{\rm ext}(\bm w)$ denote prescribed body-force and traction power.
Variations of velocity and pressure give, for admissible homogeneous velocity
tests $\bm w$ and pressure tests $q_p$,
\begin{align}
 \int_\B\tauv:\strain(\bm w)\dd A
 -\int_\B p\,\nabla\cdot\bm w\dd A&=\ell_{\rm ext}(\bm w),
 \label{eq:stokes-weak}\\
 \int_\B q_p\nabla\cdot\bm u\dd A&=0.
\end{align}
Their strong form is $-\nabla\cdot(\tauv-p\bm 1)=\bm b$ and
$\nabla\cdot\bm u=0$, with the chosen velocity/traction conditions.

Use the resistance-minus-drive sign convention
\begin{equation}
 r_V=\qcoh+r_f(V;\sigma_n,\Theta)+\eta^{\rm dpr}V
              -\mathcal T_\chi[\tauv].
\end{equation}
For one-direction friction with sticking,
\begin{equation}
 \boxed{V\geq0,\qquad r_V\geq0,\qquad Vr_V=0.}
 \label{eq:slip-kkt}
\end{equation}
At sliding points, $\mathcal T_\chi[\tauv]=\qcoh+\mu\sigma_n+
\eta^{\rm dpr}V$. At a sticking point, the equality need not hold.
In the native discrete equations the fault drive is assembled by the adjoint
of the bulk slip-strain map, not by sampling $\tau_s$ at the fault center.

\subsection{Normal stress and the limits of a scalar potential}
\label{sec:normal}

For physical total stress $\bm\sigma=\tauv-p\bm 1$, compression-positive normal
traction is $\sigma_n=p-\bm e^n\cdot\tauv\bm e^n$.
The source-audited restricted fixture uses prescribed adiabatic pressure
$P_{\rm bg}=p_{\rm ad}=10000$ Pa as the complete friction pressure contribution:
\begin{equation}
 \sigma_{n,h}=\sigma_n^{\rm bg}+p_{\rm ad},\qquad\sigma_n^{\rm bg}=0.
 \label{eq:normal-discrete}
\end{equation}
The earlier reconstruction's retained deviatoric correction was incorrect for
this selected mode. The implementation evaluates cohesive/friction resistance
on sharp-fault surface quadrature with the prescribed pressure; bulk-conjugate
shear driving uses bulk quadrature and the slip-map adjoint. There is no normal
stress or Stokes-pressure feedback to friction in this mode and no such lag to
specify. The default dynamic-pressure mode instead uses
$\sigma_n^{\rm bg}+p-\bm n\cdot\tauv_2\bm n$ at its constitutive sample.
Neither mode is silently substituted for the other.

The derivatives of \cref{eq:friction-potential} are taken with $\sigma_n$
and $\Theta$ fixed. If the actual mechanical solve updates normal traction
with $\bm u$ or $p$, the slip residual gains the corresponding Jacobian terms.
Differentiating $\mu\sigma_n(\bm u,p)V$ as a joint potential with respect to
all unknowns would add reciprocal terms to bulk equilibrium that the stated
model does not contain. Therefore no global scalar mechanical potential is
claimed for the fully normal-stress-coupled friction problem. The weak residuals
remain the governing definition.

\subsection{Evolution of stress when the profile changes}

Differentiating \cref{eq:state-reconstruction} with the rate laws gives
\begin{align}
 \dot\Eb&=\strain-\chi V\Smat-\frac{\tauv}{2\eta}
                       -s^e\dot\chi\Smat,\\
 \boxed{\dot\tauv+\frac{\tauv}\lambda
    =2G(\strain-\chi V\Smat)-2Z\dot\chi\Smat.}
 \label{eq:stress-evolving}
\end{align}
The last term is elastic redistribution at fixed committed inelastic strain.
It is signed work, not a positive frictional or viscous loss. The selected
time discretization treats it as an instantaneous phase-stage jump, rather
than spreading an additional correction strain rate over the mechanical step.

\section{Deriving the phase equation from total energy}
\label{sec:phase}

\subsection{What is held fixed}

During a phase trial, fix $\Etot$, $\Ein$, $s^e$, geometry, and material
parameters. Equivalently fix $\R$ and $Z$. Neither bulk stress nor cohesive
traction is fixed:
\begin{equation}
 \delta\tauv=-2Z\,\delta\chi\Smat,
 \quad \delta\qcoh=-\frac Z{I^2}\delta I,
 \quad \delta I=\int_{\mathcal T_\xi}h'(\phi)\delta\phi\dd\zeta,
 \quad \delta\chi=\frac{h'}I\delta\phi-\frac h{I^2}\delta I.
 \label{eq:phase-variations}
\end{equation}
The three variations are
\begin{align}
 \delta E_b&=-\int_\B\frac ZG(\tauv:\Smat)\delta\chi\dd A,
 \label{eq:bulk-var}\\
 \delta E_c&=-\int_\F\frac{Z^2}{2GI^2}\delta I\dd\xi,
 \label{eq:coh-var}\\
 \delta F&=\int_\B A\,[\alpha'(\phi)\delta\phi
                   +2\ell^2\nabla\phi\cdot\nabla\delta\phi]\dd A,
 \quad A=\frac{\Gc}{c_0\ell}.
 \label{eq:fracture-var}
\end{align}
In particular, the factor in \cref{eq:bulk-var} is $Z/G=s^e$, not $Z$.
The term involving $\delta I$ in $\delta\chi$ is essential.

\subsection{Continuum strong equation for direct normalization}

Define the normalized transverse average
$\langle\tau_s\rangle_\chi=\int\chi\tau_s\dd\zeta$.
Because $Z$ and $I$ are constant along each normal line,
\begin{align}
 \delta E_b
 &=-\int_\F\frac Z{GI}\int_{\mathcal T_\xi}
     h'(\phi)\bigl(\tau_s-\langle\tau_s\rangle_\chi\bigr)
     \delta\phi\dd\zeta\dd\xi,\\
 \delta E_c
 &=-\int_\B\frac{(\qcoh)^2}{2G}h'(\phi)\delta\phi\dd A.
\end{align}
For constant $A,\ell$, integration by parts in $\delta F$ yields
\begin{equation}
 \boxed{f_\phi=A\,[\alpha'(\phi)-2\ell^2\Delta\phi]
 -h'(\phi)\left[
       \frac{(\qcoh)^2}{2G}
      +\frac{\qcoh}G(\tau_s-\langle\tau_s\rangle_\chi)
        \right].}
 \label{eq:phase-strong}
\end{equation}
The homogeneous natural phase condition is
$\nabla\phi\cdot\bm n_{\partial\B}=0$.
Equation~\eqref{eq:phase-strong} is the direct-normalization continuum
expression. It must not be substituted for the native discrete derivative of
the projected $I_h$; that derivative is given in \cref{sec:discrete-phase}.

The first bracketed driving contribution is cohesive storage release.
The second measures the benefit or cost of redistributing slip across a
nonuniform transverse bulk stress. It can have either sign.
At uniform transverse shear stress and exact normalization the bulk term
cancels. It does not cancel in a general two-dimensional field.

\subsection{Irreversibility and stability}

For the step from the accepted state $a$, define the admissible set
\begin{equation}
 K_a=\{\phi:\ \phi^a\leq\phi\leq1,\quad
             \text{all required normalizations are finite and positive}\}.
\end{equation}
The intended constrained relaxation is associated with
\begin{equation}
 \Ecal_a(\phi)=\Ecal[\R^a,Z^a,\phi].
\end{equation}
Its first-order variational inequality is
\begin{equation}
 D\Ecal_a(\phi^k)[\psi-\phi^k]\geq0
 \quad\text{for admissible variations toward }\psi\in K_a.
 \label{eq:phase-vi}
\end{equation}
For the positive energy-gradient residual $r_{\phi,i}=\partial\Ecal_a/
\partial\phi_i$, the nodal conditions are
\begin{equation}
 \begin{cases}
 r_{\phi,i}=0,&\phi_i^a<\phi_i<1,\\
 r_{\phi,i}\geq0,&\phi_i=\phi_i^a<1,\\
 r_{\phi,i}\leq0,&\phi_i=1>\phi_i^a\quad\text{if normalization permits}.
 \end{cases}
 \label{eq:phase-kkt}
\end{equation}
The opposite residual convention reverses these signs. At coincident lower
and upper bounds the nodal value is fixed and there is no free variation.

\caution{KKT stationarity and a decrease from the incoming energy do not prove
a local minimum for this nonconvex functional. A local minimum additionally
requires a nonnegative second variation on its critical cone; a negative
eigenvalue on a genuinely free set already disproves this condition. The latest
affine-side trial has two negative free-curvature modes, and feasible
perturbations reduce its energy. The current solver therefore must be described
as a constrained stationary relaxation algorithm, not as a qualified
energy-minimizing branch-selection algorithm.}

In a smooth ideal quasistatic limit, phase irreversibility is expressed by
$\dot\phi\geq0$, $f_\phi\geq0$, $f_\phi\dot\phi=0$ away from the upper bound.
This local-in-time statement does not select the branch or determine the
physical dissipation of a finite damage jump.

\subsection{Why the old maximum-history equation is different}

Since $h'=-g'/g^2$, the direct-normalization cohesive variation alone can be
written
\begin{equation}
 -\frac{(\qcoh)^2}{2G}h'
     =g'(\phi)Y,\qquad
 Y=\frac{(\qcoh)^2}{2Gg(\phi)^2}.
\end{equation}
This explains the earlier recoverable driver, but it omits the bulk-profile
variation unless the cancellation conditions hold. Replacing $Y$ by
$\mathcal H^+(t)=\max_{r\leq t}Y(r)$ changes the phase potential again.
The selected total-energy model makes neither replacement. The initialized
particle $H$ prepares the seed, and accepted diagnostic $\mathcal H^+$ may be
recorded afterward; neither is an independent evolving source in
\cref{eq:phase-vi}.

\section{The stationary one-dimensional band}
\label{sec:stationary}

This calculation links the current energy back to the analytical profile.
Take $\phi=\phi(y)$, a horizontal fault, and uniform shear stress
$\tau_{xy}=T$. Then the bulk term in \cref{eq:phase-strong} cancels and the
free AT1 phase equation is
\begin{equation}
 A(1-2\ell^2\phi_{yy})-\frac{(\qcoh)^2}{2G}h'(\phi)=0.
 \label{eq:stationary-ode}
\end{equation}
Multiplying by $\phi_y$ and using a compact support edge with
$\phi=\phi_y=0$ gives the first integral
\begin{equation}
 \boxed{\ell^2\phi_y^2=\phi-\frac{(\qcoh)^2}{2GA}h(\phi).}
 \label{eq:stationary-first}
\end{equation}
At the center, $\phi=\widehat\phi$ and $\phi_y=0$, hence
\begin{equation}
 \frac{(\qcoh)^2}{2GA}=\frac{\widehat\phi}{h(\widehat\phi)}.
 \label{eq:stationary-core}
\end{equation}
The half support width follows by quadrature:
\begin{equation}
 b_\phi=\ell\int_0^{\widehat\phi}
 \left[z-\frac{\widehat\phi}{h(\widehat\phi)}h(z)\right]^{-1/2}\dd z.
\end{equation}
This is a stationary \emph{family}; the initial profile is not supposed to
remain unchanged under arbitrary subsequent loading.

The cohesion calibration $m=2GA/c^2$ gives
\begin{equation}
 \qcoh=c\frac{1-\widehat\phi}{\sqrt{1+p\widehat\phi}}.
 \label{eq:q-profile}
\end{equation}
The link to the original degradation relation is especially transparent here.
Let $\tau_s^{\rm tr}=(2G\R):\Smat=T+Z\chi=T+\qcoh h$.
With $T-\mu P=\qcoh$ and zero radiation damping,
\begin{equation}
 T-\mu P=g(\phi)(\tau_s^{\rm tr}-\mu P).
\end{equation}
This pointwise reduction uses uniform transverse equilibrium. In a general
field the governing friction balance is the work-conjugate surface equation,
not an additional independently imposed pointwise degradation equation.

For $p=1$, $c=1000$ Pa, and $\widehat\phi=0.6$,
$\qcoh=316.227766$ Pa. With $\mu P=1000$ Pa, the steady shear stress is
$T=1316.227766$ Pa. An exactly maintained stationary Maxwell sliding state is
\begin{equation}
 Z=I\qcoh,\qquad V=\frac Z\eta,\qquad
 u_{x,y}=\frac{T+Z\chi(y)}\eta,\qquad
 2R_{xy}=\frac{T+Z\chi(y)}G.
 \label{eq:stationary-maxwell}
\end{equation}
Both stress and $Z$ then have zero time derivative; the prescribed profile
remains a stationary constrained state. This is a \emph{loaded stationary
control}. It is distinct from the unloaded native seed with $\R=Z=0$.

The localization report verifies the stationary profile and Maxwell maintenance
to high accuracy. Its loaded reference begins instead with an unloaded seed,
then evolves to $\phi_{\max}\simeq0.85645$ at 1 s. Its phase FWHM is about
$0.38675$ m while its $\chi$ FWHM is only $0.11825$ m. A wide visible phase
profile does not imply an equally wide slip distribution.

\section{Exact Maxwell integration and endpoint equilibrium}
\label{sec:time}

\subsection{Phase redistribution at zero elapsed time}

Given $\phi^k$ from the phase step, reconstruct once from the immutable anchor:
\begin{equation}
 \boxed{\tauv^*=2G\R^a-2Z^a\chi^k\Smat
   =\tauv^a-2Z^a(\chi^k-\chi^a)\Smat,\qquad
 \qcoh{}^*=\frac{Z^a}{I^k}.}
 \label{eq:jump}
\end{equation}
The elastic cohesive slip $Z^a/G$ is unchanged. Displacement and committed
inelastic strain are unchanged. No viscous or frictional dissipation is assigned
to a zero-duration stage. Its finite energy decrease is recorded separately.

\subsection{Fixed-profile constitutive step}

Over an interval of length $\Delta t$, hold $\phi^k$, geometry, $\bm u^k$, and
$V^k$ fixed. Define
\begin{equation}
 \beta=\exp(-\Delta t/\lambda),\qquad
 \kappa=\eta(1-\beta),\qquad
 \bm B^k=\strain(\bm u^k)-\chi^k V^k\Smat.
\end{equation}
Solving the linear ODEs at an elapsed time $r\in[0,\Delta t]$ yields
\begin{align}
 \tauv(r)&=e^{-r/\lambda}\tauv^*
            +2\eta(1-e^{-r/\lambda})\bm B^k,\\
 Z(r)&=e^{-r/\lambda}Z^a+\eta(1-e^{-r/\lambda})V^k.
\end{align}
Therefore
\begin{equation}
 \boxed{\tauv^k=\beta\tauv^*+2\kappa\bm B^k,
 \qquad Z^k=\beta Z^a+\kappa V^k,
 \qquad\qcoh{}^k=\frac{Z^k}{I^k}.}
 \label{eq:endpoints}
\end{equation}
These identities hold for arbitrary positive $\Delta t/\lambda$ under the
stated constant-rate/material assumptions. The approximation
$1-\beta\simeq\Delta t/\lambda$ is not used in the model.

\subsection{Which forces are equilibrated}

The selected mechanics substitutes $\tauv^k$ and $\qcoh{}^k$ into
\cref{eq:stokes-weak,eq:slip-kkt}. It does \emph{not} substitute time means.
Nevertheless, $V^k$ and $\bm u^k$ are used as constant interval rates for slip,
deformation, and constitutive integration. Thus $V^k$ is an endpoint-force-selected
interval rate, not a reconstructed instantaneous slip rate of the exact
continuous solution. With time-dependent imposed boundary velocity, the fixtures
use its exact interval mean so that imposed displacement is correct.

For comparison, the exact mean stresses on this same constitutive path are
\begin{equation}
 r_\beta=\frac{\lambda(1-\beta)}{\Delta t},\quad
 \bar\kappa=\eta(1-r_\beta),\qquad
 \overline\tauv=r_\beta\tauv^*+2\bar\kappa\bm B^k,
 \quad\overline{\qcoh}=\frac{r_\beta Z^a+\bar\kappa V^k}{I^k}.
 \label{eq:mean-stress}
\end{equation}
Equilibrating these means defines the earlier mean scheme, not the selected
endpoint scheme. Neither choice should be substituted into the other scheme's
work ledger without the associated correction.

\subsection{An endpoint stress potential under frozen friction data}

For fixed $\sigma_n,\Theta$, the endpoint residuals can be generated by the
rate potential
\begin{align}
 \Pi_{\rm ep}(\bm u,p,V)
 &=\int_\B[\kappa|\strain-\chi V\Smat|^2
       +\beta\tauv^*:(\strain-\chi V\Smat)]\dd A\nonumber\\
 &\quad+\int_\F\left[\frac{\kappa V^2}{2I}
       +\frac{\beta Z^a V}{I}\right]\dd\xi
       +\mathcal R_f(V;\sigma_n,\Theta)\nonumber\\
 &\quad-\int_\B p\nabla\cdot\bm u\dd A-\ell_{\rm ext}(\bm u).
 \label{eq:endpoint-potential}
\end{align}
Its velocity derivative returns $\tauv^k$; its slip derivative returns the
resistance-minus-drive equation. This potential has units of power and is an
\emph{algorithmic stress potential}. It is not $\Ecal$, and it is not the
time-integrated physical work. Do not derive the phase equation by varying it
in place of \cref{eq:total-energy}. The normal-stress qualification in
\cref{sec:normal} applies when friction data are updated in the coupled solve.

\section{Derivation of the native FE strain history}
\label{sec:native}

Substitute \cref{eq:flow-laws} into the inelastic-strain rate:
\begin{align}
 \dot\R
 &=\strain-\frac{\tauv}{2\eta}-\chi\frac Z\eta\Smat\nonumber\\
 &=\strain-\frac{2G\R-2Z\chi\Smat}{2\eta}
                      -\chi\frac Z\eta\Smat
 =\strain-\frac G\eta\R.
\end{align}
Thus
\begin{equation}
 \boxed{\dot\R+\lambda^{-1}\R=\strain.}
 \label{eq:R-ode}
\end{equation}
This is an algebraic consequence of the selected constitutive state definition
and matching material parameters. It does not follow from projecting particle
stress, and it does not say that actual elastic strain is independent of damage:
$\Eb=\R-(Z/G)\chi\Smat$ still depends strongly on phase and slip.

For spatially constant $\lambda$ and a compatible initial $\R$, introduce
$\bm d_R$ with $\R=\mathcal D\bm d_R$, where $\mathcal D$ is the chosen
linear strain operator. Taking
\begin{equation}
 \dot{\bm d}_R+\lambda^{-1}\bm d_R=\bm u
\end{equation}
and applying $\mathcal D$ gives \cref{eq:R-ode}. The unloaded seed admits
$\bm d_R(0)=0$. Since $\bm u_h$ and $\bm d_{R,h}$ share the native velocity basis,
the exact constant-rate recurrence is
\begin{equation}
 \boxed{\bm d_{R,h}^k=\beta\bm d_{R,h}^a+\frac\kappa G\bm u_h^k,
 \qquad\R_h^k=\beta\R_h^a+\frac\kappa G\strain_h(\bm u_h^k).}
 \label{eq:FE-recurrence}
\end{equation}
This is a coefficient-wise update; it requires no projection of the strain
tensor into continuous Q1 fields. Derivatives of a continuous velocity FE field
need not be continuous across elements. Such derivative jumps are allowed by
this representation and are not removed by the exact history update.

Substituting \cref{eq:FE-recurrence,eq:endpoints} back into the stress gives
\begin{align}
 2G\R_h^k-2Z_h^k\chi_h^k\Smat
 &=\beta[2G\R_h^a-2Z_h^a\chi_h^k\Smat]
      +2\kappa[\strain_h(\bm u_h^k)-\chi_h^k V_h^k\Smat]\nonumber\\
 &=\beta\tauv_h^*+2\kappa\bm B_h^k.
 \label{eq:FE-Maxwell-equivalence}
\end{align}
This identity holds at every evaluation point using the same FE operator and
surface extension. It proves exact reproduction of the \emph{discrete}
Maxwell recurrence, not spatial accuracy of its stress field.

\subsection{Symmetric versus deviatoric gradients in the native report}
\label{sec:trace}

The historical native reports used raw symmetric gradients. The source audit
in \cref{sec:plane-strain} distinguishes them from the 3D plane-strain deviator
at nonzero discrete divergence. The selected correction uses $\mathcal D$
in history, stress, energy and mechanics, with the reduced 3D metric. It leaves
exact one-dimensional shear invariant. Weak Q2/Q1 incompressibility is not a
reason to omit this distinction from pointwise energy or normal traction.
A coefficient-wise recurrence remains exact because $\mathcal D$ is linear.

\subsection{What the FE history does and does not eliminate}

Particles no longer own the mechanical stress history in this option. Their
published stresses are evaluations of \cref{eq:state-reconstruction} after
acceptance and never project back into $\bm d_R$. Surface $Z$ remains a direct
surface history. Particle-domain fracture integration, material interfaces,
normalization, and geometry associations remain. The generic particle
mapping still serves diagnostics and the legacy particle-owned option.
The localization task bypassed three redundant reads; it did not remove the
generic projector or change normalization.

Spatially variable $\lambda$ prevents commuting $\mathcal D$ through
$\lambda^{-1}\bm d_R$ without additional terms. Unequal constituent parameters,
arbitrary incompatible initial strains, moving geometry, rotation, and
advection likewise require a new derivation. The current formula must not be
extended to these cases just by changing input parameters.

\section{The actual discrete phase functional and its derivatives}
\label{sec:discrete-phase}

\subsection{Evaluation and consistent normalization}

Let $\vect\phi$ be Q1 phase coefficients. Define:
\begin{itemize}
 \item $B$: evaluation of phase at bulk mechanical/energy quadrature points;
 \item $T$: evaluation of surface Q1 coefficients at the associated bulk points;
 \item $N$: evaluation of surface Q1 coefficients at surface quadrature points;
 \item $W_b,W_s$: positive diagonal bulk and surface quadrature weights;
 \item $i(\vect\phi)$: direct normal integrals at the surface quadrature abscissae.
\end{itemize}
For fixed geometry and quadrature,
\begin{equation}
 M=N^TW_sN,\quad P=M^{-1}N^TW_s,\quad
 c=P\,i(\vect\phi),\quad I_b=Tc,\quad I_s=Nc.
 \label{eq:projection}
\end{equation}
Here $c$ is a vector of normalization coefficients, not material cohesion.
Let $J_i=D_{\vect\phi}i$. Then
\begin{equation}
 \delta c=P J_i\delta\vect\phi,\qquad
 P^T=W_sNM^{-T}.
 \label{eq:projection-adjoint}
\end{equation}
The superscript $T$ on $P$ denotes its algebraic adjoint, not a replacement
normalization method. $M$ and $P$ are fixed, while $i$, $J_i$, and their second
derivatives depend on the supplied trial phase.

\subsection{Discrete energy}

Let $Z_b=T\vect Z^a$, $Z_s=N\vect Z^a$ and $h_b=h(B\vect\phi)$.
With the accepted $\R_b^a$ evaluated directly from its FE coefficients,
\begin{equation}
 \chi_b=\frac{h_b}{I_b},\qquad
 \tauv_b^*=2G\R_b^a-2Z_b\chi_b\Smat_b.
\end{equation}
All component-wise divisions here and below are at the corresponding quadrature
points. The actual discrete functional is
\begin{equation}
 \Ecal_{a,h}(\vect\phi)=
 \sum_b w_b\frac{\tauv_b^*:\tauv_b^*}{4G}
 +\sum_j w_{s,j}\frac{Z_{s,j}^2}{2G I_{s,j}}+F_h(\vect\phi),
 \label{eq:discrete-energy}
\end{equation}
with the native CPDI fracture measure
\begin{equation}
 F_h=\sum_p V_p A_p
       [\alpha(w_p\vect\phi)+\ell_p^2|D_p\vect\phi|^2].
 \label{eq:CPDI}
\end{equation}
The fixed row $w_p$ averages phase over a particle domain; $D_p$ evaluates its
domain-averaged phase gradient. $V_p$ is its physical volume/area. This expression
is not silently replaced by nodal or ordinary FE fracture quadrature.

\subsection{Exact gradient and the normalization adjoint}

Define the bulk dual load $v$ and the dual load $t$ on normalization coefficients:
\begin{align}
 v_b&=-w_b\frac{Z_b}G(\tauv_b^*:\Smat_b),\\
 t&=-T^T\diag(h_b/I_b^2)v
          -N^TW_s\frac{Z_s^2}{2G I_s^2}.
\end{align}
The positive energy-gradient residual is
\begin{equation}
 \boxed{r_\phi=F_h'
        +B^T\diag(h_b'/I_b)v+J_i^TP^Tt.}
 \label{eq:discrete-gradient}
\end{equation}
The first bulk term differentiates the numerator of $\chi$. Both denominator
effects---bulk and cohesive---pass through $P^T$ and $J_i^T$. Freezing $I_h$
inside a phase derivative, omitting $P^T$, or inserting a generic particle-to-fault
projection gives a different equation.
For AT1 with constant parameters,
\begin{equation}
 F_h'=\sum_p V_p A\,[w_p^T+2\ell^2D_p^TD_p\vect\phi],
 \qquad F_h''=2A\ell^2\sum_pV_pD_p^TD_p.
\end{equation}

\subsection{Exact second variation}

For any phase direction $d$, let $\delta\phi_b=Bd$ and
$\delta^2 c=P D^2i[d,d]$. Then
\begin{align}
 \delta\chi_b&=\frac{h_b'\delta\phi_b}{I_b}
                    -\frac{h_b\delta I_b}{I_b^2},\\
 \delta^2\chi_b&=\frac{h_b''(\delta\phi_b)^2}{I_b}
 -\frac{2h_b'\delta\phi_b\delta I_b}{I_b^2}
 -\frac{h_b\delta^2I_b}{I_b^2}
 +\frac{2h_b(\delta I_b)^2}{I_b^3}.
\end{align}
The directional Hessian is
\begin{align}
 D^2\Ecal_{a,h}[d,d]
 &=D^2F_h[d,d]\nonumber\\
 &+\sum_b w_b\left[
 \frac{2Z_b^2}{G}(\Smat_b:\Smat_b)(\delta\chi_b)^2
 -\frac{Z_b}G(\tauv_b^*:\Smat_b)\delta^2\chi_b\right]\nonumber\\
 &+\sum_j w_{s,j}\frac{Z_{s,j}^2}{G}
 \left[\frac{(\delta I_{s,j})^2}{I_{s,j}^3}
             -\frac{\delta^2I_{s,j}}{2I_{s,j}^2}\right].
 \label{eq:discrete-hessian}
\end{align}
The terms containing second derivatives of $\chi$ and $I$ make this Hessian
potentially indefinite. A positive search metric may help solve the equations,
but it does not turn the underlying energy into a convex function.

\section{Discrete endpoint mechanical system}
\label{sec:discrete-mechanics}

Let $E\vect u$ evaluate the selected FE strain operator at bulk quadrature.
Let $L(\vect\phi)$ spread surface slip coefficients into the same bulk measure:
\begin{equation}
 (L\vect V)_b=\chi_b(T\vect V)_b\Smat_b.
 \label{eq:spread-map}
\end{equation}
The bulk and surface endpoint fields are
\begin{equation}
 \vect\tau^k=\beta\vect\tau^*+2\kappa(E\vect u-L\vect V),
 \qquad \vect q_s^k=\frac{\beta N\vect Z^a+\kappa N\vect V}{I_s}.
\end{equation}
Tensor inner products are understood in $W_b$. If $D$ is the discrete divergence
matrix with pressure integration absorbed, the residuals have the form
\begin{align}
 r_u&=E^TW_b\vect\tau^k-D^T\vect p-\vect f,\\
 r_p&=-D\vect u,\\
 r_V&=-L^TW_b\vect\tau^k
       +N^TW_s[\vect q_s^k+\vect r_f+\eta^{\rm dpr}N\vect V].
 \label{eq:mechanical-residuals}
\end{align}
The $V$ equations are imposed as equalities on sliding/free degrees of freedom
and with the selected sticking complementarity at contacts. Prescribed slip
degrees of freedom instead have reactions.

For frozen normal stress and state, the principal coupled Jacobian blocks are
\begin{align}
 K_{uu}&=2\kappa E^TW_bE,&
 K_{uV}&=-2\kappa E^TW_bL,\\
 K_{Vu}&=-2\kappa L^TW_bE,&
 K_{VV}&=2\kappa L^TW_bL+
 N^TW_s\diag\!\left(\frac\kappa{I_s}
              +\partial_Vr_f+\eta^{\rm dpr}\right)N.
\end{align}
Normal-stress/state feedback adds its actual chain-rule terms to the slip
rows. No symmetry claim applies to those additions. In particular,
work conjugacy requires the \emph{same} $L$ and $L^T$ and the same bulk
quadrature, whether or not friction makes the full tangent nonsymmetric.

The restricted native quadrature splits bulk cells at surface-Q1 intersections,
uses Q3 on the resulting tangential intervals and Q3 normally, and uses the
same selected points for the associated mechanical blocks and bulk energy.
Correct history ownership and exact normalization cannot compensate for an
underresolved velocity space.

\section{One accepted native step: the algorithm}
\label{sec:algorithm}

The following is a mathematical specification of the selected order.
\begin{enumerate}[label=\textbf{\arabic*.},itemsep=6pt]
 \item \textbf{Initialize once.} Prepare the nominal phase band using the native
 prescribed-$H$ relaxation on $[0,1]$. Commit its result as pre-existing damage.
 Set $\bm d_R^0=0$, $Z^0=0$, stress and accumulated slip to zero. Compute the
 consistent normalization. The unloaded seed is protected by the phase obstacle;
 it is not declared an unconstrained minimum of the unloaded total energy.

 \item \textbf{Snapshot the accepted state.} Fix $\phi^a$, $\bm d_R^a$, $Z^a$,
 mesh, particle domains, geometry, and associations. Compute $\beta,\kappa$.
 Use $\phi^a$ as the immutable lower bound for every phase trial. Generic
 extrapolated \texttt{current\_linearization\_point}, \texttt{old\_solution},
 and \texttt{old\_old\_solution} do not replace this owner or its phase anchor.

 \item \textbf{Solve the phase stationarity problem.} For each supplied trial
 $\vect\phi$, recompute direct normal integrals, their consistent projection,
 $\chi$, reconstructed $\tauv^*$, the actual energy, gradient, and tangent.
 Enforce admissibility and \cref{eq:phase-kkt}. Use the native constrained
 search and unchanged convergence target. The new exact free-set Newton
 polish is a small numerical correction near a root, not a new physical law
 or a minimum-selection criterion. Record $\phi^k$ only as staged data.

 \item \textbf{Record instantaneous redistribution.} Form $\tauv^*$ using
 \cref{eq:jump}; compute the bulk and cohesive storage jumps, fracture
 expenditure, and phase-release quantity in \cref{sec:ledger}. Do not
 dissipate these jumps through an invented viscous time interval.

 \item \textbf{Solve endpoint mechanics.} Hold $\phi^k,I_h^k,\chi_h^k$
 fixed and solve \cref{eq:mechanical-residuals} for the interval rates
 $\bm u_h^k,V_h^k$ and pressure. All nonlinear candidates reconstruct the
 same $\tauv^*,Z^a$; iterating the solve must not repeatedly advance history.
 Use the configured normal-stress/friction convention and slip constraints.

 \item \textbf{Stage history updates.} Apply
 \[
 \bm d_R^k=\beta\bm d_R^a+(\kappa/G)\bm u^k,\qquad
 Z^k=\beta Z^a+\kappa V^k,\qquad
 s^k=s^a+\Delta t V^k.
 \]
 Reconstruct stress and cohesive traction directly from these staged fields.
 When displacement is recorded, $\bm d^k=\bm d^a+\Delta t\bm u^k$;
 it is a different field from $\bm d_R^k$.

 \item \textbf{Validate and publish once.} Check phase inequalities,
 normalization, the actual selected residuals, FE recurrence, Maxwell
 consistency, and separate energy quantities. Prepare derived particle
 stress and diagnostic histories. Publish the staged FE/surface/phase state
 only at terminal acceptance. Failed phase or mechanical trials must not
 become an accepted owner. This local ownership discipline does not by
 itself qualify full simulator rollback, restart, or transported histories.
\end{enumerate}

The final accepted mechanical state has not yet driven a phase solve. Solving
phase again with that new $\R^k,Z^k$ generally produces another change.
This is the stated one-step splitting lag, not evidence that the recorded
phase solve failed its own earlier-anchor residual. The localization report's
additional final-anchor phase change is about $0.00928$ in the rich reference.

\section{Energy accounting: what closes and what does not}
\label{sec:ledger}

\subsection{Continuous smooth identity}

With compatible bulk/slip equilibrium, \cref{eq:energy-rate-preflow} gives
\begin{equation}
 \dot\Ecal+\Dcal_b+\Dcal_c+\Dcal_f+\Dcal_r
  =\mathcal P_{\rm ext}+D_\phi\Ecal[\dot\phi],
\end{equation}
where $\Dcal_f=\int_\F r_f V\dd\xi$ and
$\Dcal_r=\int_\F\eta^{\rm dpr}V^2\dd\xi$.
Smooth phase complementarity makes the last term zero. This is not a proof
that a finite-step sequential scheme, or a finite phase jump, has zero
unaccounted energy.

\subsection{Phase relaxation release}

At fixed accepted $\R^a,Z^a$, let
\begin{align}
 J_b&=E_b(\R^a,Z^a,\phi^k)-E_b(\R^a,Z^a,\phi^a),\\
 J_c&=\int_\F\frac{(Z^a)^2}{2G}
             \left(\frac1{I^k}-\frac1{I^a}\right)\dd\xi,\\
 \Delta F&=F(\phi^k)-F(\phi^a),\\
 \boxed{D_{\rm relax}=-(J_b+J_c+\Delta F).}
 \label{eq:Drelax}
\end{align}
A decreasing-energy phase step has $D_{\rm relax}\geq0$ up to its arithmetic
allowance. It is the finite-drop slack \emph{after} fracture expenditure has
already been counted. It is neither an additional fracture energy nor a
prescribed material damage viscosity. Its physical interpretation or vanishing
limit is an unresolved qualification when it is appreciable.

\subsection{Exact fixed-profile constituent identities}

Using the same spatial quadratures as the equations, integrate the constitutive
paths exactly or by independently verified temporal quadrature. Define
\begin{align}
 W_{\rm int}&=\int_0^{\Delta t}\int_\B\tauv(r):\strain\dd A\dd r,\\
 W_{\rm slip}&=\int_0^{\Delta t}\int_\B\tau_s(r)\chi^k V^k\dd A\dd r,\\
 W_q&=\int_0^{\Delta t}\int_\F\qcoh(r)V^k\dd\xi\dd r.
\end{align}
The local ODEs guarantee
\begin{equation}
 \Delta E_b^{\rm mech}+D_b=W_{\rm int}-W_{\rm slip},\qquad
 \Delta E_c^{\rm mech}+D_c=W_q.
 \label{eq:constituent-ledgers}
\end{equation}
Here $D_b=\int\Dcal_b\dd r$ and $D_c=\int\Dcal_c\dd r$.
With constant interval resistance, $D_f=\Delta t\int r_f^k V^k\dd\xi$;
if resistance is instead reevaluated on the Maxwell path, its actual time
integral must be used and that choice must be stated. This matters when the
normal-stress correction is nonzero. The formulas below are valid for either
explicitly specified convention.

\subsection{The extra endpoint mechanical discrepancy}

Let $W_{\rm boundary}$ include all actual external work, including reactions
on prescribed velocity boundaries and any other prescribed-rate work.
Define
\begin{align}
 E_{\rm exchange}&=W_{\rm slip}-W_q-D_f-D_r,\\
 W_{\rm free}&=W_{\rm int}-W_{\rm boundary},\\
 \boxed{E_{\rm mech}=E_{\rm exchange}-W_{\rm free}.}
 \label{eq:mechanical-error}
\end{align}
Combining the phase and mechanical stages gives the exact accounting identity
\begin{equation}
 \boxed{W_{\rm boundary}
  -\Delta(E_b+E_c+F)-D_b-D_c-D_f-D_r
  =D_{\rm relax}+E_{\rm mech}.}
 \label{eq:full-ledger}
\end{equation}
Endpoint force balance sets the appropriate end-of-interval residuals to zero.
It does not force their integrals over the Maxwell path to vanish. Therefore
$E_{\rm mech}$ generally remains nonzero. Equating all bulk internal work to
external boundary work would hide $W_{\rm free}$; adding $E_{\rm mech}$ to
phase dissipation would hide a different error. Report signed and accumulated
absolute discrepancies separately.

Exact Maxwell integration, a small nonlinear residual, and accurate $I_h$
are each valuable but none implies that \cref{eq:full-ledger} has a negligible
right-hand side. The endpoint comparison shows decreasing discrepancies under
time refinement; the latest native localization tests are at one timestep and
do not establish native time convergence.

\section{Interpreting the latest localization report}
\label{sec:report}

The physical setup uses $G=10^6$ Pa, $\eta=10^8$ Pa s, $\Gc=26.6666666667$
J/m$^2$, $\ell=0.15625$ m, $m=128$, and $\Delta t=1/64$ s. These are dimensional
verification parameters, not nondimensionalized rock properties.

\subsection{Boundary conditions explain the gross delocalization}

The all-affine control prescribes $u_x=0.256t(y-0.5)$ and $u_y=0$ on all
four sides. This forbids the non-affine $u_x(y)$ needed to accommodate a
localized slip band at the side walls. It reaches nearly full-height phase
and slip-distribution widths at 1 s.

The compatible control prescribes both components only on top and bottom;
on vertical sides it prescribes $u_y=0$ and leaves $u_x$ free. The associated
natural normal traction fixes Stokes pressure, so the old arbitrary pressure
normalization is disabled. The background friction pressure remains prescribed.
This control retains a localized band and is constant along the fault to roundoff.

\begin{center}\small
\begin{tabular}{L{.30\textwidth}rrrr}
\toprule
Case at 1 s & $\phi_{\max}$ & $\phi$ FWHM & $\chi$ FWHM & $\qcoh$ center\\
 & & m & m & Pa\\\midrule
Affine sides, height 1 & .65836 & 1.00000 & 1.00000 & 143.13\\
Compatible, height 1 & .82144 & .38795 & .13305 & 141.58\\
Compatible, height 2 & .88743 & .44950 & .13670 & 97.17\\
Rich 1D, height 1 & .85645 & .38675 & .11825 & 91.04\\
\bottomrule
\end{tabular}\end{center}

Doubling the height while keeping the same local imposed shear-rate function
also doubles the imposed boundary velocity difference and final relative
displacement. That comparison changes physical compliance/loading distance;
it is not a pure boundary-artifact test.

\subsection{Correct native equations are not yet accurate native stresses}

The compatible native run matches the \emph{same coarse} Q1-phase/Q2-velocity
reference to about $8.72\times10^{-9}$ in phase. Against the rich reference,
however, the maximum common-time phase error is $0.04642$ and the cohesive
traction error is $51.83$ Pa. Bulk shear stress oscillates by tens of kPa around
a mean of order 1 kPa. Enriching only the velocity representation substantially
improves onset and phase response. The exact FE history recurrence has removed
one transfer inconsistency; it has not removed approximation error in
$2G\R_h-2Z_h\chi_h\Smat$.

\subsection{Stationarity is a separate issue from stability}

The former failed phase entry was not limited by the measured arithmetic floor.
An exact free-set Newton correction reduced its residual without changing the
target. But its free Hessian has two negative modes. The implemented polish
repairs equation solving; it does not repair the energy landscape or select a
stable branch. The appropriate next mathematical work is therefore native
spatial accuracy and a declared stability/branch-selection policy, followed by
coupled time/energy qualification. Changing $G$, weakening normalization
accuracy, or adding a damage cap would not address the identified causes.

\section{Current equations versus superseded interpretations}

\begin{longtable}{L{.29\textwidth}L{.61\textwidth}}
\toprule Earlier interpretation & Current restricted P2 definition\\\midrule\endhead
Maximum single-step work increment drives phase & Full instantaneous state
energy $E_b+E_c+F$ is differentiated. No timestep-sized work maximum drives phase.\\
Recoverable peak $\mathcal H^+$ alone drives phase & Retained only as diagnostic
after seeded initialization; direct irreversibility is the phase constraint.\\
Only fault cohesive energy changes with $\phi$ & Bulk stress changes through
$\chi$ at fixed $\R,Z$; both storage variations enter the phase derivative.\\
Mean forces are the selected mechanical equations & Endpoint forces are selected;
mean stresses remain necessary for exact work diagnostics.\\
Particle stress is authoritative everywhere & Optional native FE state owns
$\R$; legacy particle ownership remains the default outside this option.\\
$\bm d_R$ is displacement or ordinary viscous strain & It encodes
$\Etot-\Ein$ and has its own exponential recurrence.\\
A KKT pass proves phase energy minimization & It proves first-order constrained
stationarity only; negative curvature has actually been observed.\\
Accurate normalization implies accurate traction & Velocity-space stress errors
can remain very large with accurate $I_h$ and exact history updates.\\
\bottomrule
\end{longtable}

\appendix
\section{Analytic normal integral for Q1 phase, \texorpdfstring{$p=1$}{p=1}}

On a normal segment of length $L$, let phase vary linearly from $a$ to $b$.
Set $u=1-a$, $v=1-b$, and $z=(v-u)/u$. Since
$h=m(2/(1-\phi)^2-3/(1-\phi)+1)$,
\begin{equation}
 \int_0^L h(\phi(\zeta))\dd\zeta
 =mL\left[\frac2{uv}-\frac{3\log(1+z)}{uz}+1\right].
 \label{eq:analytic-I}
\end{equation}
The $z=0$ limit is $Lh(a)$. A stable implementation uses
\texttt{log1p} and its continuous small-$z$ limit. Differentiating this same
integral with respect to its endpoint values supplies $J_i$ and $D^2i$.
This removes normal quadrature error for the represented Q1 field, not phase
interpolation error or surface projection error. Nonintegrable saturated
segments remain invalid.

\section{Exact time kernels for work and viscous losses}

Let $x=\Delta t/\lambda$ and
\begin{align}
 J_{ee}&=\frac\lambda2(1-e^{-2x}),\\
 J_{e1}&=\frac\lambda2(1-e^{-x})^2,\\
 J_{11}&=\Delta t-2\lambda(1-e^{-x})+J_{ee}.
\end{align}
For $A(r)=e^{-r/\lambda}A_0+\eta(1-e^{-r/\lambda})b$,
\begin{equation}
 \int_0^{\Delta t}|A(r)|^2\dd r
 =J_{ee}|A_0|^2+2\eta J_{e1}A_0:b+\eta^2J_{11}|b|^2.
\end{equation}
Use $A_0=\tauv^*$, $b=2\bm B^k$ for bulk stress and $A_0=Z^a$, $b=V^k$
for cohesive history; divide by $2\eta$ and $\eta I^k$, respectively, to obtain
the viscous losses. Integrated stresses are
\begin{align}
 \int_0^{\Delta t}\tauv\dd r
   &=\lambda(1-\beta)\tauv^*
      +2\eta[\Delta t-\lambda(1-\beta)]\bm B^k,\\
 \int_0^{\Delta t}Z\dd r
   &=\lambda(1-\beta)Z^a
      +\eta[\Delta t-\lambda(1-\beta)]V^k.
\end{align}
Stable \texttt{expm1} or series evaluations near zero prevent cancellation;
using a numerical series to evaluate these kernels does not change the
general-Maxwell constitutive law into a small-step physical approximation.

\section{Direct one-dimensional endpoint reference}

Let $r(y)=2R_{xy}(y)$, domain height $L_y$, mean imposed shear rate
$\dot\gamma=(u_x^{\rm top}-u_x^{\rm bottom})/L_y$, and a newly fixed profile.
Set $\tau_0(y)=G r^a(y)-Z^a\chi(y)$. A uniform endpoint stress $T$ obeys
\begin{equation}
 T=\beta\tau_0(y)+\kappa[u_{x,y}(y)-\chi(y)V].
\end{equation}
Integrating across the height gives
\begin{equation}
 T=\beta\overline\tau_0+\kappa(\dot\gamma-C_QV/L_y),
 \qquad C_Q=\sum_q w_q\chi_q,
 \qquad \overline\tau_0=\frac1{L_y}\sum_q w_q\tau_{0,q}.
\end{equation}
Retain the sampled $C_Q$ rather than assuming it is exactly one. With constant
friction resistance $\mu P$, zero radiation damping, and a positive-slip branch,
the work-conjugate equation is
$C_QT=(\beta Z^a+\kappa V)/I+\mu P$. Therefore
\begin{equation}
 V=\max\left\{0,
 \frac{C_Q(\beta\overline\tau_0+\kappa\dot\gamma)
                -\beta Z^a/I-\mu P}
      {\kappa(C_Q^2/L_y+1/I)}\right\},\qquad
 u_{x,y}=\frac{T-\beta\tau_0}{\kappa}+\chi V.
\end{equation}
The accepted continuous reference is $r^k(y)=[T+Z^k\chi(y)]/G$.
Its nonpolynomial profile is retained directly, rather than forcing it into
coarse Q2 derivatives. This is why it can distinguish constitutive/history
correctness from velocity-space approximation error.

\section{Where a future RSF law would enter}

The current controls have no evolving $\Theta$. A later RSF test would replace
$r_f$ in the same slip residual by the selected rate-and-state resistance and
add its state equation, for example the aging-law update
\begin{equation}
 \dot\Theta=1-\frac{V\Theta}{D_c},\qquad
 \Theta^k=\frac{\Theta^a+\Delta t}{1+\Delta t V^k/D_c}
\end{equation}
if backward Euler is the selected state integrator. The chosen regularized
friction function, state lag/coupling, normal-traction evaluation, reverse-slip
handling, and Jacobian must then be specified and tested. This is an extension
location, not evidence that native FE history has already qualified BP3, RSF
instability, or a general energy-minimization formulation for stateful friction.

\section{Source hierarchy and remaining documentation gaps}

The most recent localization report controls the interpretation of the current
native option. The native FE and endpoint reports supply the state ownership
and time scheme. The older total-energy algorithm note supplies the energy
derivative derivation but described a mean-force prototype; its choice of mean
equilibrium is superseded here. The original manuscript supplies notation and
historical context; its incremental-work phase driver is not restored.

This note is self-contained and requires no external figures or bibliography
files. To compile, run \texttt{pdflatex phase\_field\_native\_fe\_derivation.tex}
twice. The plane-strain and friction-pressure audit above supersedes the
original attachment-only qualifications concerning those source definitions.
The design authority remains current\_design.md and specification.tex; this
note derives the scientific embedding and records the bounded source correction.
General dynamic-normal-stress friction work, spatial convergence, minimum
selection and native coupled temporal accuracy remain distinct qualifications.

\subsection{Bounded post-audit accuracy and stability evidence}
The corrected optional Newton matrix agrees with the explicit embedded operator
to relative $4.92\times10^{-14}$; a compatible shear control is unchanged.
Two compatible native controls with Q2/Q1 and Q3/Q2 mechanics finish to
$0.625$ s with unchanged phase bounds and KKT tolerances, no measured healing,
and an exact saved FE coefficient recurrence. Q3 enrichment improves common-time
phase error against the rich one-dimensional reference from $0.03643$ to
$0.01312$. This is not spatial convergence: the retained three-point normal
quadrature gives nearly constant Q3 stress at its samples, while independent
12/24-point frozen evaluation gives $409$ Pa transverse stress RMS and bulk
storage $0.875859$ J/m instead of the native $0.792491$ J/m. Normalization
representation error is negligible for these along-fault-uniform fields.

At sampled loaded states, the free phase Hessians have positive minimum
eigenvalues and feasible small-mode perturbations increase energy. This does
not test the complete constrained critical cone or all states; the old affine
negative-curvature evidence remains separate. In particular, a phase-independent
trace-storage correction at a fixed accepted anchor cannot remove those modes.
No quadrature or phase-branch correction is inserted on the strength of these
checks. See \path{propagation/P2_plane_strain_report.md} for settings, profiles,
the source-consumer inventory, separate energy errors, and actual source and
specification diffs. Full spatial/stability/time qualification and P3 remain
deferred.

The subsequent coupled-quadrature assessment uses six normal points throughout
the native calculation, selected by inexpensive four/six/twelve-point controls
and independent 24/48-point evaluation. It retains the physical settings and
consistent normalization/CPDI integrations. Q3 reduces the common-time phase
error by about $2.8$ and transverse stress RMS by $7.2$ at $2.11$ times native
runtime. A normal-refined $16\times32$ Q2 comparator meets the declared restricted
phase/traction/slip targets at similar cost; both still have substantial stress
and bulk-storage spatial error. New sampled curvature checks use the proper
incoming anchors. These results supersede the three-point accuracy assessment,
whose fields and failed evidence remain preserved; they do not establish full
stability, time accuracy or P3 readiness. See
\path{propagation/P2_bulk_quadrature_report.md} for definitions, costs and limits.

\clearpage
\begin{thebibliography}{9}
\small
\bibitem{localization}
Codex project report, \emph{P2 native localization and numerical diagnosis},
9 October 2026. Supplied file: \path{P2_native_localization_report.md}.

\bibitem{native}
Codex project report, \emph{Restricted native P2: accepted FE strain}.
Supplied file: \path{P2_native_FE_strain_report(1).md}.

\bibitem{endpoint}
Codex project report, \emph{Evolving endpoint/mean comparison},
8 October 2026. Supplied file: \path{P2_total_energy_endpoint_comparison.md}.

\bibitem{algorithm}
Codex derivation, \emph{One total-energy candidate step},
8 October 2026. Supplied file: \path{P2_total_energy_algorithm.md}.

\bibitem{manuscript}
Yimin Jin, \emph{Phase-Field Model for Rate-State-Friction Law in Stokes System}.
Supplied manuscript: \path{pf_rsf(20261007-203939).tex}.
\end{thebibliography}

\section{Optional native consistent-endpoint reduced action}
This default-off solution algorithm retains the energy and constitutive laws.
For every trial, fix the incoming $(\bm d_R^a,\vect Z^a,\vect\phi^a)$;
solve momentum, incompressibility and contact at the trial phase and reconstruct
$\bm d_R^b=\beta\bm d_R^a+(\kappa/G)\bm u$ and
$\vect Z^b=\beta\vect Z^a+\kappa\vect V$. The obstacle is the previous physical
phase, not a predictor or a Newton iterate. Publication occurs once.

For a direction $d$, evaluate $\delta c_I=P J_i d$ (normalization coefficients),
$\delta I_b=T\delta c_I$ and $\delta I_s=N\delta c_I$, with the actual projection
in \cref{eq:projection}. Define
\[
 \delta\chi=h'Bd/I_b-h\delta I_b/I_b^2,\qquad
 \delta\tau_{\rm partial}=-2Z_b\delta\chi\Smat.
\]
In the native residual convention the mechanical perturbation system is
\[
 \begin{pmatrix}A&-B\\G&-K_V\end{pmatrix}
 \binom{\delta(\bm u,p)}{\delta\vect V}
 =\binom{2E^TW_b(Z_b\delta\chi\Smat)}{-f_\Gamma},
\]
with zero continuity forcing, homogeneous perturbation boundary conditions,
and zero increments at sticking/prescribed surface degrees of freedom. Here
\[
 f_\Gamma=T^TW_b\{\delta\chi[\tau:\Smat-2Z_b\chi(\Smat:\Smat)]\}
              +N^TW_s Z_s\delta I_s/I_s^2.
\]
Retain the full native velocity and pressure spaces, all free slip coefficients,
and tangential modes. With $\delta Z=\kappa\delta V$,
\[
 \delta\tau=2\kappa D_3(\operatorname{symgrad}\delta\bm u)
             -2(\delta Z\chi+Z\delta\chi)\Smat.
\]
Let $c_b=-Z_b(\tau:\Smat)/G$ (an unweighted bulk coefficient) and
$\delta c_b=-[\delta Z_b(\tau:\Smat)+Z_b(\delta\tau:\Smat)]/G$.
The directional phase gradient is
\begin{align*}
 \delta r_\phi={}&K_Fd+
 B^TW_b\left\{\delta c_b h'/I_b+
 c_b[h''Bd/I_b-h'\delta I_b/I_b^2]\right\}\\
 &+J_i^TP^T\delta t+(D J_i[d])^TP^Tt,\\
 \delta t={}&-T^TW_b\left\{\delta c_b h/I_b^2+c_bh'Bd/I_b^2
                      -2c_bh\delta I_b/I_b^3\right\}\\
 &-N^TW_s\left\{Z_s\delta Z_s/(GI_s^2)
                         -Z_s^2\delta I_s/(GI_s^3)\right\}.
\end{align*}
The analytic normal-ray Hessians implement the final adjoint term without a
dense native phase Hessian. $K_F$ remains the actual native CPDI fracture block.
Only the homogeneous bulk matrix and preconditioners are reused within a step;
phase-dependent $B,G,K_V$ are refreshed at a new trial.
The restricted native option stores the existing exact sparse assembly of $B$
to avoid repeated quadrature actions; the integration and constraints are the
same. It also reuses the existing exact sparse $G$ assembly, with the same
normal-filter guard, constitutive coefficients and remote-point averaging.
This stores the current linear map, rather than replacing it by $B^T$ or
changing the pressure convention. Rebuild at every surface linearization;
physical pressure scaling and homogeneous constraints remain in the condensed
solver. A homogeneous sparse velocity factor serves as a within-step
preconditioner, without eliminating pressure, slip or tangential modes.

Fresh uncondensed residual corrections refine the mechanical tangent on this
same operator. Each solve retains the configured Stokes linear tolerance.
The change in the resulting elastic phase action is compared against the
configured phase linear tolerance times its norm; unresolved accuracy fails.
This is an action-uncertainty estimate, checked against finite differences and
independent tighter-inner repeats. It does not replace the true outer KKT,
actual-energy line search, physical obstacle, or terminal history checks.


\paragraph{Obstacle-aware Newton search (numerical correction).}
At each nonstationary iterate $x$, predict a search active set with the existing
positive fracture-plus-mass metric $P$, using
\[
 \widehat s=\arg\min_{s\ge\phi^a-x}
       \{r_\phi(x)^Ts+\tfrac12s^TPs\}.
\]
This inexpensive auxiliary solve is not a phase equation. With its proposed
active set $\mathcal A$ and free set $\mathcal F$, solve the actual reduced
Newton system
\[
 d_{\mathcal A}=\phi^a_{\mathcal A}-x_{\mathcal A},\qquad
 (\mathsf H_{\rm red})_{\mathcal F\mathcal F}d_{\mathcal F}
 =-(r_\phi)_{\mathcal F}
  -(\mathsf H_{\rm red})_{\mathcal F\mathcal A}d_{\mathcal A}.
\]
Compute the active-column term with the same linearized mechanical/normalization
pullback, keeping the current trial and incoming accepted histories fixed.
Use the configured relative linear tolerance on this free right-hand side.
The initial feasible metric step and a descent fallback remain search guesses;
publication still requires the original true KKT, actual-energy line search,
physical bounds and curvature checks. This avoids admitting one normal layer
per nonlinear iteration after a clipped step. It changes neither the governing
energy nor its Hessian, and does not advance history or physical time.

\paragraph{Intermediate numerical accuracy and profile-translation inverse.}
The optional adaptive search changes intermediate Newton stopping, not the
physical admissible set or final KKT. For projected defects $d_k,d_0$, the
unchanged entry target $t_\phi$ and configured linear accuracy $\eta_{\min}$,
use
\[
 \eta_k=\min\{.01,\max\{\eta_{\min},.1\sqrt{d_k/\max(d_0,t_\phi)}\}\},
 \qquad t_{\rm lin}=\max\{\eta_k\|b_{\mathcal F}\|_2,.1t_\phi\}.
\]
The actual energy line search and independent original CPDI KKT remain unchanged.
Mechanical tangent-action refinement retains $\eta_{\min}$.

One targeted preconditioning block uses $-\partial_y\phi$ with independent
tangential nodal amplitudes, masked to the proposed free set. Early adaptive
searches use one global column; when $\eta_k<.01$, expand to the 17 columns of
this fixed fixture. This addresses measured tangential late directions without
restricting any native unknown. Evaluate every image with the actual coupled
Hessian and diagonalize only its small symmetric coarse matrix. Resolved positive
eigenpairs give $Q=\sum_j u_ju_j^T/\lambda_j$; numerical positivity uses the
existing action-accuracy scale. The balanced inverse is
\[
 C_{\rm prof}=(I-Q\mathsf H_{\rm red})C(I-\mathsf H_{\rm red}Q)+Q,
\]
with the left factor implemented as the transpose of the right for the inexact
action. Omitting negative coarse modes modifies the preconditioner only. The
physical Hessian, actual negative-curvature checks, accepted reference and
publication rules remain fixed. No centering or additional dissipation is implied.

After a resolved saddle departure, restart the intermediate forcing reference
from the new defect if larger. The physical entry target is never restarted;
accepted histories and the obstacle remain fixed.

Near a soft branch, sum CPDI fracture increments before adding their fixed
reference. Extended accumulation of the same storage and friction integrands
reduces summation noise in the true-energy line search; gradients, Hessians,
quadrature and acceptance allowances remain unchanged.

\end{document}
